% GENERATED FILE. Edit canonical lessons, then run npm run generate:books.
% canonical-source-hash: fnv1a64:7ed0c2f245d017af
% canonical-lessons: MW-C140-undro, MW-C140-makhi, MW-C140-machchhar, MW-C140-kiri, MW-C140-bichhu

\chapter{Mouse, Fly, Mosquito, Ant, Scorpion}
\label{ch:mw-mouse-fly-mosquito-ant-scorpion}
\hlchaptermodality{140}

\begin{chapteropening}
\textbf{By the end of this chapter:} \emph{I can say a mouse, a fly, a mosquito, an ant and a scorpion.}

Complete the last lesson of chapter 140: I can say a mouse, a fly, a mosquito, an ant and a scorpion.
\end{chapteropening}

\section[\texorpdfstring{\mw{उंदरो} (undro)}{उंदरो (undro)}]{\mw{उंदरो} (undro) — a mouse}

\label{lesson:MW-C140-undro}



\begin{quote}
Before the new one: say the Marwadi for a koel, then the Marwadi for a kite (the bird).
\end{quote}

\subsection*{You'll want to know: \mw{उंदरो}}
\textbf{\mw{उंदरो}} — \emph{undro} — \textquotedblleft{}a mouse\textquotedblright{}.

\begin{cousinweb}
\textbf{\mw{उंदरो}} (\emph{undro}) is a tadbhava, a word inherited from Sanskrit and reshaped by sound change. It comes from \emph{undura}, \textquotedblleft{}rat, mouse\textquotedblright{}. CDIAL gives no Marwari form in this entry (CDIAL 2095); the match is through Gujarati \emph{ũdar} and Hindi \emph{ũdur}. CDIAL takes the Sanskrit word itself from an Austroasiatic language.
\end{cousinweb}

\begin{grammarlens}[title={What you've built}]
One more word for this chapter.
\end{grammarlens}

\subsection*{Your turn}
\begin{itemize}
\raggedright
  \item \emph{Say it:} \emph{undro}
  \item \emph{Say it:} \emph{undro}, once more
  \item \emph{Say it:} say \emph{undro}
  \item \emph{Recall:} say the Marwadi for a koel, then the Marwadi for a kite (the bird), then say \emph{undro} again
\end{itemize}

\begin{tcolorbox}[breakable,colback=teal!4,colframe=teal!35!black,title={Before you move on}]
What does \mw{उंदरो} mean? (\textquotedblleft{}A mouse\textquotedblright{}.) Say it once more.
\end{tcolorbox}
\section[\texorpdfstring{\mw{माखी} (mākhī)}{माखी (mākhī)}]{\mw{माखी} (mākhī) — a fly}

\label{lesson:MW-C140-makhi}



\begin{quote}
Before the new one: say the Marwadi for a kite (the bird), then the Marwadi for a mouse.
\end{quote}

\subsection*{You'll want to know: \mw{माखी}}
\textbf{\mw{माखी}} — \emph{mākhī} — \textquotedblleft{}a fly\textquotedblright{}.

\begin{cousinweb}
\textbf{\mw{माखी}} (\emph{mākhī}) is a tadbhava, a word inherited from Sanskrit and reshaped by sound change. It comes from \emph{makṣikā}, \textquotedblleft{}fly, bee\textquotedblright{}. CDIAL gives no Marwari form in this entry (CDIAL 9696); the match is through Hindi \emph{mākhī} and Gujarati \emph{mākhī}.
\end{cousinweb}

\begin{grammarlens}[title={What you've built}]
One more word for this chapter.
\end{grammarlens}

\subsection*{Your turn}
\begin{itemize}
\raggedright
  \item \emph{Say it:} \emph{mākhī}
  \item \emph{Say it:} \emph{mākhī}, once more
  \item \emph{Say it:} say \emph{mākhī}
  \item \emph{Recall:} say the Marwadi for a kite (the bird), then the Marwadi for a mouse, then say \emph{mākhī} again
\end{itemize}

\begin{tcolorbox}[breakable,colback=teal!4,colframe=teal!35!black,title={Before you move on}]
What does \mw{माखी} mean? (\textquotedblleft{}A fly\textquotedblright{}.) Say it once more.
\end{tcolorbox}
\section[\texorpdfstring{\mw{मच्छर} (machchhar)}{मच्छर (machchhar)}]{\mw{मच्छर} (machchhar) — a mosquito}

\label{lesson:MW-C140-machchhar}



\begin{quote}
Before the new one: say the Marwadi for a mouse, then the Marwadi for a fly.
\end{quote}

\subsection*{You'll want to know: \mw{मच्छर}}
\textbf{\mw{मच्छर}} — \emph{machchhar} — \textquotedblleft{}a mosquito\textquotedblright{}.

\begin{cousinweb}
\textbf{\mw{मच्छर}} (\emph{machchhar}) is a tadbhava, a word inherited from Sanskrit and reshaped by sound change. It comes from \emph{matsara}, \textquotedblleft{}mosquito, fly\textquotedblright{}. CDIAL gives no Marwari form in this entry (CDIAL 9757); the match is through Hindi \emph{macchar}. CDIAL notes the possibility that the Sanskrit word was made back from Prakrit \emph{macchara}, but finds no support for it.
\end{cousinweb}

\begin{grammarlens}[title={What you've built}]
One more word for this chapter.
\end{grammarlens}

\subsection*{Your turn}
\begin{itemize}
\raggedright
  \item \emph{Say it:} \emph{machchhar}
  \item \emph{Say it:} \emph{machchhar}, once more
  \item \emph{Say it:} say \emph{machchhar}
  \item \emph{Recall:} say the Marwadi for a mouse, then the Marwadi for a fly, then say \emph{machchhar} again
\end{itemize}

\begin{tcolorbox}[breakable,colback=teal!4,colframe=teal!35!black,title={Before you move on}]
What does \mw{मच्छर} mean? (\textquotedblleft{}A mosquito\textquotedblright{}.) Say it once more.
\end{tcolorbox}
\section[\texorpdfstring{\mw{कीड़ी} (kīṛī)}{कीड़ी (kīṛī)}]{\mw{कीड़ी} (kīṛī) — an ant}

\label{lesson:MW-C140-kiri}



\begin{quote}
Before the new one: say the Marwadi for a fly, then the Marwadi for a mosquito.
\end{quote}

\subsection*{You'll want to know: \mw{कीड़ी}}
\textbf{\mw{कीड़ी}} — \emph{kīṛī} — \textquotedblleft{}an ant\textquotedblright{}.

\begin{cousinweb}
\textbf{\mw{कीड़ी}} (\emph{kīṛī}) is a tadbhava, a word inherited from Sanskrit and reshaped by sound change. It comes from \emph{kīṭa}, \textquotedblleft{}insect, worm\textquotedblright{}. CDIAL gives no Marwari form in this entry (CDIAL 3193); the match is through Hindi \emph{kīṛī} and Gujarati \emph{kīṛī}. Hindi and Gujarati \emph{kīṛī} also mean \textquotedblleft{}ant\textquotedblright{}.
\end{cousinweb}

\begin{grammarlens}[title={What you've built}]
One more word for this chapter.
\end{grammarlens}

\subsection*{Your turn}
\begin{itemize}
\raggedright
  \item \emph{Say it:} \emph{kīṛī}
  \item \emph{Say it:} \emph{kīṛī}, once more
  \item \emph{Say it:} say \emph{kīṛī}
  \item \emph{Recall:} say the Marwadi for a fly, then the Marwadi for a mosquito, then say \emph{kīṛī} again
\end{itemize}

\begin{tcolorbox}[breakable,colback=teal!4,colframe=teal!35!black,title={Before you move on}]
What does \mw{कीड़ी} mean? (\textquotedblleft{}An ant\textquotedblright{}.) Say it once more.
\end{tcolorbox}
\section[\texorpdfstring{\mw{बीछू} (bīchhū)}{बीछू (bīchhū)}]{\mw{बीछू} (bīchhū) — a scorpion}

\label{lesson:MW-C140-bichhu}



\begin{quote}
Before the new one: say the Marwadi for a mosquito, then the Marwadi for an ant.
\end{quote}

\subsection*{You'll want to know: \mw{बीछू}}
\textbf{\mw{बीछू}} — \emph{bīchhū} — \textquotedblleft{}a scorpion\textquotedblright{}.

\begin{cousinweb}
\textbf{\mw{बीछू}} (\emph{bīchhū}) is a tadbhava, a word inherited from Sanskrit and reshaped by sound change. It comes from \emph{vṛścika}, \textquotedblleft{}scorpion\textquotedblright{}. CDIAL gives no Marwari form in this entry (CDIAL 12081); the match is through Hindi \emph{bīchū} and Gujarati \emph{vīchu}. CDIAL suggests that the many forms of the word come from a taboo on naming the scorpion.
\end{cousinweb}

\begin{grammarlens}[title={What you've built}]
One more word for this chapter.
\end{grammarlens}

\subsection*{Your turn}
\begin{itemize}
\raggedright
  \item \emph{Say it:} \emph{bīchhū}
  \item \emph{Say it:} \emph{bīchhū}, once more
  \item \emph{Say it:} say \emph{bīchhū}
  \item \emph{Recall:} say the Marwadi for a mosquito, then the Marwadi for an ant, then say \emph{bīchhū} again
\end{itemize}

\begin{tcolorbox}[breakable,colback=teal!4,colframe=teal!35!black,title={Before you move on}]
What does \mw{बीछू} mean? (\textquotedblleft{}A scorpion\textquotedblright{}.) Say it once more.
\end{tcolorbox}
